\documentclass[12pt]{amsart}
\input{../../../../../Headers/TestPackageHeader}

\renewcommand{\headrulewidth}{0pt} % remove lines as well

\numberwithin{equation}{section}

\DeclareGraphicsRule{.tif}{png}{.png}{`convert #1 `dirname #1`/`basename #1 .tif`.png}

\textheight = 10.8 in
\topmargin = -1.0in
\textwidth = 7.7in
\oddsidemargin = -0.7in
\evensidemargin = -0.7in

\begin{document}

\input{../../../../../Headers/TestDefinitionHeader.tex}

$\quad$

\vspace{1.0cm}

\begin{center}
\LARGE {MA 301} \\

\Large{Test 1} \\

\large{Fall - 2011} \\

\vspace{0.5cm}

\setlength{\extrarowheight}{0.4cm}
\begin{tabular}{|c|c|c|}
\hline
Problem & points & out of  \\
\hline
\hline
1 & & 10 \\
\hline
2 & & 10 \\
\hline
3 & & 10 \\
\hline
4 & & 10 \\
\hline
5 & & 10 \\
\hline
\hline
Score & & 50\\
\hline
\end{tabular}
\vspace{0.5cm}

%--------TABLE---------
	\setlength{\extrarowheight}{0.2cm}
	\begin{tabular}{rcc}
		Name: & \underline{$\qquad \qquad \qquad \qquad \quad$}& \\
		&&$\hspace{1.4cm}$\\
		HR: & \underline{$\qquad \qquad \qquad \qquad \quad$}& \\
		&&\\
		Signature: & \underline{$\qquad \qquad \qquad \qquad \quad$}& \\
	\end{tabular}
%--------------------

\vspace{0.5cm}

	\setlength{\extrarowheight}{0.2cm}
	\begin{tabular}{| p{12cm} |}
		\hline
		\multicolumn{1}{|c|}{{\bf READ BEFORE YOU START}} \\
		\hline
		 * Show all work.\\
		 * Change your calculator mode to RADIANS.\\
		 * Sloppy solutions may be penalized.\\
		 * Use calculators ONLY for computations.\\
		 * You must write up every problem in a clear and \\
		 $\,\,$ organized way.\\
		 * If you write solutions on a page different than the \\
		 $\,\,$ question, then clearly state where it is located.\\
		\hline
	\end{tabular}

\end{center}

\vspace{1.0cm}

\thispagestyle{empty}
\newpage
\thispagestyle{empty}

{\bf Problem 1:} Given $\text{\bf F} = [2,2y+z,y+1]$

(a) Show that $\text{\bf F}$ is conservative in two different ways

	\begin{tabular}{p{10cm}|p{10cm}}
		\multicolumn{1}{c|}{method 1}&\multicolumn{1}{c}{method 2}\\
		&\\&\\&\\&\\&\\&\\&\\&\\&\\&\\&\\&\\&\\&\\&\\&\\&\\&\\&\\&\\&\\&\\&\\&\\&\\&\\&\\&\\&\\&\\&\\&\\&\\&\\&\\&\\&\\
	\end{tabular}

\vspace{0.4cm}

(b) Evaluate $\displaystyle{\int \limits_{C} \text{\bf F} \cdot  d\text{\bf r}}$ when $C$ is a straight line from $(-1,4,2)$ to $(0,0,-6)$

\thispagestyle{empty}
\newpage
\thispagestyle{empty}

{\bf Problem 2:} Let $\text{\bf F} = [x+y,2xy]$ and $C$ be the boundary of the region bounded by $y=x^2$ and $y=x+2$ in the $x,y$-plane. \\
(a) Compute \,\, $\D{\oint \limits_{C} \text{\bf F}\cdot d\text{\bf r}}$ \,\, using green's theorem.\\
(b) \UL{Setup} the integral you would use to compute \,\, $\D{\oint \limits_{C} \text{\bf F}\cdot d\text{\bf r}}$ \,\, directly \,\,\,({\em don't leave answer as a dot product})

\thispagestyle{empty}
\newpage
\thispagestyle{empty}

{\bf Problem 3:} Let $\text{\bf F} = [x,y,z]$ and $S$ be the sphere $x^2+y^2+z^2 = 25$.  \,\,Compute \,$\D{\iint \limits_{S} \text{\bf F}\cdot \text{\bf n} \, dA}$

\vspace{0.2cm}

\thispagestyle{empty}
\newpage
\thispagestyle{empty}

{\bf Problem 4:} Compute the flux of water (water density $= 1 \, ton/m^3$) through the parabolic cylinder $S:\, y=x^2$ where $0 \le x \le 2$ and $0 \le z \le 1$ if the velocity flow of water is given by $\text{\bf v} = [3z^2, 6, 6xz]$ (units: {\em m/s}).

\thispagestyle{empty}
\newpage
\thispagestyle{empty}

{\bf Problem 5:} Answer the following conceptual questions

\vspace{0.1cm}

(a) \quad Suppose $\text{{\bf F}} = [1,1,1]$.  Is it true that \,\,$\D{\iint \limits_{S} \text{\bf F}\cdot \text{\bf n} \, dA = 0}$ \,\,for any closed surface $S$ which is a boundary of a solid?  Justify your answer.

\vspace{4.0cm}

(b) \quad Describe the difference between Green's Theorem and Stokes' Theorem.

\vspace{4.0cm}

(c) \quad Briefly describe the meaning of flux.

\vspace{4.0cm}

(d) Let $\text{\bf F}$ vector field such that curl\,$\text{\bf F} =$ {\bf 0}.  Match each expression in the left column with the most appropriate term in the right column.

\begin{center}
%\setlength{\extrarowheight}{0.8cm}
\begin{tabular}{cp{5cm}c}
	$\D{\int \limits_{C} \text{\bf F}\cdot d\text{\bf r}}$ & & conservative\\
	&&\\
	$\D{\text{\bf F}\cdot d\text{\bf r}}$ & & zero\\
	&&\\
	$\D{\oint \limits_{C} \text{\bf F}\cdot d\text{\bf r}}$ & & path independent\\
	&&\\
	$\D{\text{\bf F}}$ & & exact\\
\end{tabular}
\end{center}

\vspace{0.4cm}

(e) Suppose we know to represent a hill as a surface given by the function $f(x,y)$.  If we place a ball anywhere on the hill, what direction would you expect it to roll when we let it go?

 
\thispagestyle{empty}
\newpage
\thispagestyle{empty}

\end{document}